\section{El contrato del modelo de mundo: estado, acción y futuro predecible}

Esta clase no trata el ``world model'' como un generador de video más grande, sino que primero establece qué contrato debe cumplir: dado el estado actual y las acciones ejecutables, el modelo debe predecir cómo evoluciona el entorno, permitiendo que la predicción sea suficiente para razonar, planificar y controlar. Los dos ponentes responden a este contrato por dos rutas complementarias: Causal-JEPA añade sesgos estructurales de interacción de objetos a las tareas de predicción, mientras que LeWorldModel comprime el entrenamiento extremo a extremo en un objetivo estable y conciso.

\subsection{De la predicción de un paso a la transferencia de estado controlada}

La predicción autoregresiva estándar se escribe simplemente como $x_t\mapsto x_{t+1}$, pero el futuro en entornos reales también depende de las acciones tomadas por el agente y los participantes externos. Por ello, el modelo de mundo mínimo adoptado en esta clase es la transferencia de estado controlada: el modelo debe saber no solo ``qué hay ahora'', sino también ``qué cambiará tras aplicar una acción''. Al leer la primera figura, se deben entender las tres columnas como tres criterios de aceptación independientes, no como una propiedad que una red satisfaga automáticamente simultáneamente.

\lecturefigure{slide-003.jpg}{Los tres contratos del modelo de mundo: representación de estado, leyes de transferencia y respuesta a acciones}{CS25 V6 oficial deck, p.~3}

\[
f_\theta:\mathcal{S}\times\mathcal{A}\rightarrow\mathcal{S},
\qquad
\hat{s}_{t+1}=f_\theta(s_t,a_t).
\]
Donde $\mathcal{S}$ es el espacio de estados, $\mathcal{A}$ es el espacio de acciones, $s_t$ es el estado en el tiempo $t$, $a_t$ es la acción aplicada en ese momento, $f_\theta$ es el modelo de transferencia parametrizado por $\theta$, y $\hat{s}_{t+1}$ es el siguiente estado predicho por el modelo. Si el entorno contiene aleatoriedad no observable, una escritura más honesta sería la distribución condicional $p_\theta(s_{t+1}\mid s_{\le t},a_{\le t})$, en lugar de asumir un futuro único y determinado.

\teachervoice{00:01:52--00:03:41，Hazel denomina intuitivamente al world model como simulator y pide a los lectores que recuerden siempre tres preguntas de diseño: ¿la representación conserva la información importante actual?, ¿el modelo de transferencia ha aprendido las reglas del entorno?, y ¿las acciones cambian el futuro de la manera correcta?}

\begin{importantbox}{Un modelo de mundo no es un solo nombre de red}
Un modelo de mundo funcional contiene al menos tres capas de interfaz: el \textbf{observation encoder} convierte píxeles en estados, el \textbf{dynamics predictor} proyecta el futuro y el \textbf{decision interface} conecta las predicciones con la búsqueda de acciones o el aprendizaje de políticas. Solo mostrar un video realista demuestra que el generador produce alguna secuencia visual, pero no prueba que su estado sea controlable, que las acciones tengan significado semántico correcto o que sus predicciones sean suficientes para planificar.
\end{importantbox}

\subsection{Por qué esta clase no busca por ahora reconstruir cada píxel}

Los modelos de mundo generativos pueden producir futuros altamente realistas, pero el requisito de verosimilitud a nivel de píxel exige que el modelo explique simultáneamente la estructura relevante para la tarea y una gran cantidad de detalles impredecibles, como texturas, ruido de iluminación o movimiento aleatorio en el fondo. Los ponentes no niegan Sora, GAIA, Genie ni la generación interactiva en 3D; más bien, estrechan el problema a: si el objetivo es el control y la inferencia, ¿se puede aprender directamente un \textit{latent state} que retenga solo los factores predecibles y decidibles?

\lecturefigure{slide-004.jpg}{Modelos de mundo generativos fuera del alcance de esta clase: video, conducción y entornos interactivos}{CS25 V6 oficial deck, p.~4}

En la siguiente página, las diferencias se dibujan en la ubicación de la señal de supervisión. Las arquitecturas generativas envían el resultado de la predicción al espacio de píxeles para compararlo directamente con el futuro real $y$; las arquitecturas \textit{joint-embedding} codifican también el objetivo como $s_y$, requiriendo solo que el \textit{latent} predicho sea compatible con el \textit{latent} objetivo. El núcleo no es simplemente eliminar el decodificador, sino permitir que el espacio de representación ignore los detalles que no afectan a la tarea entre múltiples futuros posibles.

\lecturefigure{slide-005.jpg}{Las ubicaciones de supervisión difieren entre predicción generativa y predicción \textit{joint-embedding}}{CS25 V6 oficial deck, p.~5}

\[
\mathcal{L}_{\text{pixel}}=\ell\!\left(D(P(E(x),a)),y\right),
\qquad
\mathcal{L}_{\text{latent}}=\ell\!\left(P(E(x),a),\operatorname{sg}(E(y))\right).
\]
Donde $E$ es el codificador, $P$ es el predictor y $D$ es el decodificador; $x$ es el contexto, $a$ es la acción o condición auxiliar, $y$ es la observación futura, $\ell$ es una distancia o verosimilitud logarítmica negativa, y $\operatorname{sg}$ indica stop-gradient. La primera ecuación debe reconstruir el dominio objetivo; la segunda solo requiere alineamiento en el espacio de representación, por lo que se pueden retener selectivamente los factores predecibles.

\begin{warningbox}{\textit{latent} no es automáticamente igual a \textit{semantic}}
Mover el \textit{loss} al espacio \textit{latent} solo cambia la interfaz de optimización; si el codificador aprende constantes, texturas de fondo o atajos del conjunto de entrenamiento, el \textit{latent} aún puede ser inútil. El \textit{object masking} posterior y SIGReg resuelven respectivamente ``¿de qué depende la predicción?'' y ``¿cómo evitar que la representación colapse?'', y no pueden sustituirse mutuamente.
\end{warningbox}

\subsection{Perspectiva basada en energía de JEPA y colapso de representación}

JEPA se puede entender mediante un \textit{energy-based model} (modelo de energía): las combinaciones contexto--futuro compatibles deben obtener baja energía, mientras que las incompatibles deben obtener alta energía. Esta perspectiva permite que un contexto corresponda a múltiples futuros razonables sin requerir promediarlos en píxeles borrosos. Sin embargo, si el codificador mapea todas las entradas al mismo vector, cualquier predicción será ``igual'' al objetivo, resultando en una pérdida de cero sin haber aprendido nada del mundo.

\lecturefigure{slide-006.jpg}{JEPA: aprendiendo compatibilidad y terreno energético en el espacio de representación}{CS25 V6 oficial deck, p.~6}

\[
\mathcal{E}_\theta(x,y)=\left\|P_\theta(E_\theta(x))-E_\theta(y)\right\|_2^2,
\qquad
\mathcal{E}_\theta(x,y^+) < \mathcal{E}_\theta(x,y^-).
\]
Donde $\mathcal{E}_\theta$ es la energía de compatibilidad, $y^+$ es el futuro compatible con $x$, $y^-$ es el futuro incompatible, y $E_\theta$ con $P_\theta$ son respectivamente el representador y el predictor. Baja energía significa ``puede ser un futuro razonable'', lo cual no es equivalente a una normalización de probabilidades de píxeles.

\teachervoice{00:06:13--00:08:16，el docente enfatiza que \textit{anti-collapse} se divide aproximadamente en dos familias: contrastiva y regularización. V-JEPA usa un codificador objetivo con EMA, stop-gradient y enmascaramiento de tareas; su propósito no es decorar el flujo de entrenamiento, sino excluir soluciones triviales como representaciones constantes.}

\lecturefigure{slide-007.jpg}{V-JEPA: enmascaramiento espacio-temporal, context encoder, target encoder y predictor}{CS25 V6 oficial deck, p.~7}

\begin{knowledgebox}{Uso por primera vez: EMA y stop-gradient}
El \textbf{EMA (Exponential Moving Average) target encoder} actualiza la rama objetivo con el promedio móvil exponencial de los parámetros del codificador en línea, haciendo que los cambios en el objetivo sean más lentos; el \textbf{stop-gradient} impide que la rama objetivo sea arrastrada directamente por el error de predicción actual. Ambos suelen estabilizar juntos el entrenamiento auto-supervisado, pero son heurísticas de optimización y no equivalen a explicar por qué la representación necesariamente no colapsa.
\end{knowledgebox}

\subsection{De aprendizaje de representación a modelos de mundo condicionados a acción}

V-JEPA 2 separa el aprendizaje de representaciones de video a gran escala del post-entrenamiento \textit{action-conditioned} con datos de interacción limitados: primero se aprende un estado visual de las observaciones, luego se congela o reutiliza la representación para entrenar el predictor condicionado a acción. DINO-WM pregunta aún más: dado que DINOv2 ya puede proporcionar características visuales fuertes, ¿es necesario volver a aprender el codificador, bastando con entrenar la dinámica sobre estas características?

\lecturefigure{slide-008.jpg}{V-JEPA 2: de pre-entrenamiento con video enmascarado a post-entrenamiento condicionado a acción}{CS25 V6 oficial deck, p.~8}

\lecturefigure{slide-009.jpg}{DINO-WM: congelar DINOv2 y predecir autoregresivamente el futuro sobre patch latent}{CS25 V6 oficial deck, p.~9}

\teachervoice{00:08:16--00:10:17，Hazel pide a los lectores distinguir dos fases: el pre-entrenamiento a gran escala libre de acción de V-JEPA 2 aprende primero la representación, mientras que V-JEPA 2-AC y DINO-WM conectan la acción y la propiocepción al predictor. La clase cuestiona posteriormente si la \textit{patch representation} es adecuada para expresar interacciones entre objetos, sin negar la capacidad de planificación de estas líneas base.}

\begin{knowledgebox}{Términos clave: cuatro conceptos fáciles de confundir}
\begin{tabularx}{\textwidth}{L{0.20\textwidth} X}
\toprule
Término & Significado en esta clase \\
\midrule
representation model & Comprime las observaciones en un estado útil para la tarea; no necesariamente sabe las consecuencias de las acciones. \\
dynamics model & Predice la evolución del estado dado el estado y la acción; puede construirse sobre representaciones congeladas. \\
world model & Combinación de representación, dinámica e interfaz de decisión; la capacidad de planificación depende de la calidad conjunta de los tres. \\
JEPA & Marco de entrenamiento que realiza predicciones en el espacio de representación; tanto el codificador como el predictor pueden implementarse completamente con Transformers. \\
\bottomrule
\end{tabularx}
\end{knowledgebox}

\subsection{Resumen de la sección}

El contrato mínimo del modelo de mundo es la transferencia de estado controlada; el valor de JEPA es mover el objetivo de predicción al espacio de representación, evitando que el modelo reconstruya todos los píxeles; su riesgo fundamental es que el \textit{latent} pueda colapsar o aprender atajos incorrectos. V-JEPA, V-JEPA 2 y DINO-WM ofrecen tres combinaciones estables de representación y dinámica, mientras que la pregunta del próximo capítulo es: ¿estas composiciones deben consistir en parches o en objetos interactivos?
